\chapter{Phương pháp đề xuất}
\label{ch:phuongphap}

\section{Mô hình hệ thống và yêu cầu thiết kế}
\label{sec:systemmodel}

\paragraph{Reframe bài toán.} Đối chiếu với hệ thống LBS thực tế
(Chương~\ref{ch:lienquan}) cho thấy dịch vụ không cần vị trí \emph{chính xác} mà
cần \emph{mức granularity tối thiểu đủ cho một mục đích}. Apple và Android đều
phơi bày các tầng độ chính xác theo query, xử lý ``significant location''
on-device và làm mờ điểm đầu/cuối hành trình; quy định (GDPR Điều~5; ngưỡng
``precise geolocation'' 1.750--1.850 ft của luật bang Mỹ; các lệnh FTC với
Kochava/X-Mode/InMarket) áp đặt purpose limitation, data minimization và loại trừ
vị trí nhạy cảm. Tuy nhiên thứ được triển khai thực tế chỉ là \emph{lượng tử hóa
heuristic không đảm bảo} (snap-to-grid, làm tròn) hoặc \emph{DP hình thức chỉ ở
phía máy chủ trên dữ liệu tổng hợp}. Luận văn vì thế phát biểu lại bài toán: từ
``làm nhiễu mỗi điểm trong bán kính QoS $\delta$ toàn cục'' thành \textbf{phát ra
mức tối thiểu theo mục đích, on-device, trên lưới đường, với đảm bảo nhận thức
endpoint và tương quan, đánh giá theo threat model rút từ các tấn công thật}.

\paragraph{Mô hình hệ thống.} Kiến trúc \emph{đích} đặt cơ chế \emph{client-side}
(mô hình local), tại tầng vị trí của thiết bị --- tiền lệ đã triển khai cho hàng
tỷ máy là Android \texttt{LocationFudger} và iOS approximate location, cùng
Location Guard/LP-Doctor cho planar Laplace. Đây là nơi đặt duy nhất mà đảm bảo
giữ được với \emph{mọi} bên downstream (máy chủ LBS, ad exchange, data broker,
trát tòa). \emph{Lưu ý phạm vi:} artifact trong luận văn là một \emph{reference
mechanism} + simulator Flask phía server, \emph{chưa} phải một hiện thực mobile
on-device; các vấn đề tích hợp OS, secure cache, TTL/rotation định danh là ngoài
phạm vi (Chương~\ref{ch:tongket}). Cơ chế được biết: (a)~vị trí thật hiện tại;
(b)~\emph{đồ thị đường công khai lưu trên thiết bị} ($\sim$5,9MB cho bbox thí
nghiệm, nhỏ hơn nhiều so với bản đồ offline 200--950MB mà điện thoại vẫn chứa ---
điều kiện hóa trên nó không tốn privacy, là post-processing); (c)~\emph{các điểm
đã CÔNG BỐ trước đó của chính nó} (không phải vị trí thật quá khứ). Adversary
thấy: điểm công bố on-road cộng metadata tầng truyền (timestamp, nội dung query,
định danh có thể liên kết) --- mô hình hóa là LBS honest-but-curious hoặc broker
giữ toàn bộ chuỗi công bố dưới một định danh.

\paragraph{Yêu cầu thiết kế và phạm vi trung thực.} Bảng~\ref{tab:requirements}
liệt kê các yêu cầu rút từ industry/pháp lý và harm chúng đối ứng. Cơ chế đề xuất
(REM/T-REM/SM-REM) làm được R1,R3,R4,R5,R6 và một phần R2,R7; \emph{không} tự làm
được R8 (xoay định danh --- cần policy phía định danh) và R9--R12 (consent, giới
hạn mục đích/lưu trữ, minh bạch --- thuộc governance phía máy chủ hoặc tầng UX).
Đây là ranh giới đóng góp cần nêu thẳng: một cơ chế nhiễu tọa độ không thể thay
thế kiểm soát định danh và quản trị dữ liệu.

\begin{table}[ht]
\centering
\caption{Yêu cầu thiết kế R1--R7 (phần cơ chế giải quyết) và harm đối ứng
(S1--S6, xem Chương~\ref{ch:lienquan}).}
\label{tab:requirements}
\small
\begin{tabular}{clll}
\toprule
\# & Yêu cầu & Nguồn & Harm \\
\midrule
R1 & Perturbation client-side/on-device & Apple; AOSP; GDPR & tiền đề \\
R2 & Granularity theo mục đích (chọn từ context công khai) & GDPR 5(1)(c); Micinski'13 & S1,S6 \\
R3 & Bất khả phân biệt hình thức vị trí gần & W3C; Geo-I & S1,S5 \\
R4 & Output on-road hợp lệ & LocationFudger; REM & S2 \\
R5 & Kháng tương quan/velocity (điều kiện trên điểm đã công bố) & Xiao--Xiong; T-REM & S3 \\
R6 & Ổn định khi báo lặp (anti-averaging) & W3C; RAPPOR; LP-Doctor & S4 \\
R7 & Bảo vệ endpoint/stay-point hạng nhất & Apple; Strava; Dhondt'22 & S4,S6 \\
\bottomrule
\end{tabular}
\end{table}

\section{Phân tích baseline Thực tập 2}
\label{sec:phantichbaseline}

Pipeline Thực tập 2 gồm bốn bước cho mỗi điểm thật $x_t$: (i) cộng nhiễu Laplace
phẳng với bán kính bị chặn tại ngưỡng QoS $\delta$; (ii) nếu điểm giả rơi vào
tòa nhà/mặt nước thì lấy mẫu lại (tối đa $N$ lần); (iii) làm mượt liên tục theo
hai điểm giả trước, có kiểm tra lại QoS so với điểm thật; (iv) làm khớp về đỉnh
đường gần nhất nếu vẫn thỏa QoS. Hai mệnh đề sau chỉ ra các bước (i)--(iii) phá
vỡ đảm bảo hình thức. (Trong benchmark Chương~\ref{ch:thucnghiem}, baseline được
dùng là một \emph{surrogate đơn giản hoá} --- bỏ bước tìm đường thay thế và bước
reject theo tòa nhà/mặt nước --- nên nhãn ``Baseline TT2 (surrogate)''; so sánh
đối chứng với pipeline đầy đủ là hướng phát triển.)

\begin{menhde}[Chặn bán kính quanh điểm thật phá vỡ Geo-I]
\label{prop:cap}
Gọi $M_\delta$ là cơ chế cộng nhiễu với bán kính $r' = \min(r, \delta)$ quanh vị
trí thật $x$. Khi đó $M_\delta$ không thỏa $\eps$-Geo-I với bất kỳ $\eps$ hữu hạn
nào trên miền $\mathcal{X}$ có đường kính lớn hơn $2\delta$.
\end{menhde}

\begin{proof}
Support của $M_\delta(x)$ là hình tròn đóng $\bar{B}(x, \delta)$. Chọn $x, x'$
với $\dEu(x, x') > 2\delta$ thì $\bar{B}(x,\delta) \cap \bar{B}(x',\delta) =
\emptyset$. Với $S = \bar{B}(x, \delta)$ ta có $\Pr[M_\delta(x) \in S] = 1$
trong khi $\Pr[M_\delta(x') \in S] = 0$; không tồn tại hằng số hữu hạn
$e^{\eps\, \dEu(x,x')}$ chặn được tỷ số này.
\end{proof}

Cần tách ba trường hợp khác nhau (theo phản biện V-010/R2-011); chỉ trường hợp
(b) mới là lỗi của baseline:

\begin{menhde}[Điều kiện hóa trên tập cho phép công khai cố định --- giữ $2\eps$]
\label{prop:fixedaccept}
Cho tập cho phép \emph{cố định, công khai} $A$ (độc lập với input) với
$K(A\mid x)>0$ mọi $x$. Cơ chế điều kiện $K_A(z\mid x)=K(z\mid x)/K(A\mid x)$ thỏa
\[
\frac{K_A(z\mid x)}{K_A(z\mid x')}
= \frac{K(z\mid x)}{K(z\mid x')}\cdot\frac{K(A\mid x')}{K(A\mid x)}
\le e^{\eps \dEu(x,x')}\cdot e^{\eps \dEu(x,x')} = e^{2\eps \dEu(x,x')}.
\]
\end{menhde}

\begin{menhde}[Lấy mẫu lại theo vùng cấm \emph{phụ thuộc input} phá vỡ Geo-I]
\label{prop:reject}
Nếu tập cho phép $A_x$ phụ thuộc input (ví dụ ``không rơi vào tòa nhà/mặt nước
quanh $x$''), tỷ số có thể vô hạn. Phản ví dụ tối thiểu: base kernel $K$ đều trên
$\{0,1\}$ với mọi input (0-DP); đặt $A_a=\{0\}$, $A_b=\{1\}$. Sau điều kiện,
$K_A(0\mid a)=1$ nhưng $K_A(0\mid b)=0$ --- likelihood ratio vô hạn. Pipeline
Thực tập 2 dùng đúng dạng $A_x$ này (vùng cấm + kiểm QoS quanh $x$).
\end{menhde}

\noindent (Trường hợp (c) --- retry hữu hạn rồi fallback --- tạo một phân phối
khác với conditioning vô hạn và phải phân tích riêng; không dùng trong các cơ chế
đề xuất.) Cùng với Nhận xét~\ref{rmk:expbug} (bán kính nhiễu lấy sai phân phối),
các phân tích cho thấy $\eps$ danh nghĩa của baseline không phải một đảm bảo hợp
lệ; mọi so sánh ``cùng $\eps$'' vì vậy chỉ mang tính tham khảo.
Ngược lại, phép chiếu đầu ra về đỉnh đường gần nhất ở bước (iv) --- nếu bỏ điều
kiện kiểm tra QoS so với điểm thật --- là hậu xử lý độc lập dữ liệu và bảo toàn
mức riêng tư; quan sát này dẫn tới thiết kế dưới đây.

\section{REM: Cơ chế lũy thừa trên lưới đường}
\label{sec:rem}

Cho đồ thị đường $G = (V, E)$ trích từ OSM cho vùng quan tâm; tập đỉnh $V$ là
\emph{công khai}. REM chọn trực tiếp một đỉnh làm vị trí công bố:

\begin{dinhnghia}[REM]
\label{def:rem}
Với vị trí thật $x$ và ngân sách $\eps$:
\[
\Pr[\mathrm{REM}(x) = v] \;=\; \frac{\exp\!\big(-\tfrac{\eps}{2}\,\dEu(x, v)\big)}
{\sum_{u \in V} \exp\!\big(-\tfrac{\eps}{2}\,\dEu(x, u)\big)},
\qquad v \in V.
\]
\end{dinhnghia}

\begin{dinhly}[REM thỏa $\eps$-Geo-I]
\label{thm:rem}
Với mọi $x, x' \in \Real^2$ và mọi $v \in V$:
$\Pr[\mathrm{REM}(x) = v] \le e^{\eps\,\dEu(x,x')}\,\Pr[\mathrm{REM}(x') = v]$.
\end{dinhly}

\begin{proof}
Đặt $Z(x) = \sum_{u \in V} e^{-\frac{\eps}{2} \dEu(x,u)}$. Xét tỷ số
\[
\frac{\Pr[\mathrm{REM}(x) = v]}{\Pr[\mathrm{REM}(x') = v]}
= \underbrace{\frac{e^{-\frac{\eps}{2}\dEu(x,v)}}{e^{-\frac{\eps}{2}\dEu(x',v)}}}_{(A)}
\cdot \underbrace{\frac{Z(x')}{Z(x)}}_{(B)}.
\]
Với $(A)$: theo bất đẳng thức tam giác
$\dEu(x',v) - \dEu(x,v) \le \dEu(x,x')$, nên
$(A) = e^{\frac{\eps}{2}(\dEu(x',v) - \dEu(x,v))} \le e^{\frac{\eps}{2}\dEu(x,x')}$.
Với $(B)$: với mỗi $u \in V$,
$\dEu(x', u) \ge \dEu(x, u) - \dEu(x, x')$, do đó
$e^{-\frac{\eps}{2}\dEu(x',u)} \le e^{\frac{\eps}{2}\dEu(x,x')}
e^{-\frac{\eps}{2}\dEu(x,u)}$. Lấy tổng theo $u$:
$Z(x') \le e^{\frac{\eps}{2}\dEu(x,x')} Z(x)$, tức
$(B) \le e^{\frac{\eps}{2}\dEu(x,x')}$. Nhân hai chặn được
$e^{\eps\,\dEu(x,x')}$.
\end{proof}

\begin{nhanxet}[Quan hệ với GEM và điều kiện metric của guarantee]
REM chính là bản dùng khoảng cách \emph{Euclid} của Graph-Exponential Mechanism
(GEM) trong~\cite{takagi2019geo}: GEM chấm điểm bằng khoảng cách đường đi ngắn
nhất $d_s$ và chỉ thỏa $\eps$-GG-I (không có đảm bảo theo Euclid), trong khi REM
chấm điểm bằng $d_{\text{Euclid}}$ nên Định lý~\ref{thm:rem} cho $\eps$-Geo-I
\emph{Euclid} thật sự. Điểm mấu chốt: guarantee bám vào metric dùng trong số mũ ---
dùng Euclid thì phát biểu Euclid, dùng graph thì chỉ được phát biểu GG-I. Việc
miền đầu ra rời rạc/on-road không phá lập luận exponential mechanism.
\end{nhanxet}

\begin{nhanxet}[Tính thực tế theo cấu trúc]
Mọi đầu ra của REM là một đỉnh đường: các tấn công lọc theo tính hợp lý không
gian (map-matching, phát hiện điểm trên mặt nước --- cơ sở của
RAoPT~\cite{buchholz2022raopt}) không còn thông tin để khai thác.
\end{nhanxet}

\begin{nhanxet}[Tập ứng viên cố định --- không dùng cutoff]
\label{rmk:cutoff}
Hiện thực lấy tổng chuẩn hóa trên \emph{toàn bộ} tập đỉnh $V$ của đồ thị (cố định,
công khai), không dùng bất kỳ cutoff nào quanh điểm thật. Đây là điều kiện để
Định lý~\ref{thm:rem} đúng: nếu cắt support về $V \cap B(x,R)$ (phụ thuộc điểm
thật $x$), hai điểm $\eps$-sát nhau nhưng ở hai phía biên cutoff của một đỉnh $v$
cho $P[M(x){=}v]>0$ trong khi $P[M(x'){=}v]=0$ --- likelihood ratio vô hạn, phá
vỡ pure Geo-I (khối lượng đuôi nhỏ không cứu được). Chi phí full-support cho đồ
thị Bắc Kinh ($|V|\approx 78$k đỉnh) là một 2-norm vector hóa cộng một lần lấy mẫu
categorical mỗi release, dưới mili-giây --- trong ngân sách thời gian thực của LBS.
\end{nhanxet}

\section{T-REM: Bổ sung ràng buộc thời gian với trọng số công khai}
\label{sec:trem}

Nhược điểm của REM (và mọi cơ chế nhiễu độc lập theo thời gian) là hai điểm công
bố liên tiếp có thể cách nhau phi thực tế, tạo tín hiệu cho tấn công liên kết
vận tốc. T-REM khắc phục bằng trọng số khả đạt tính từ điểm \emph{đã công bố}
$z_{t-1}$ và khoảng thời gian $\Delta t$ (đều công khai):
\[
w_t(v) = \exp\!\Big(-\lambda \cdot \max\big(0,\; \dEu(z_{t-1}, v) - v_{\max}\Delta t - s\big)\Big),
\]
với $v_{\max}$ là tốc độ tối đa hợp lý, $s$ là biên nới lỏng, $\lambda$ là hệ số
phạt. Cơ chế trở thành
\[
\Pr[\mathrm{TREM}(x_t) = v \mid z_{t-1}]
= \frac{w_t(v)\, e^{-\frac{\eps}{2}\dEu(x_t, v)}}
{\sum_{u \in V} w_t(u)\, e^{-\frac{\eps}{2}\dEu(x_t, u)}}.
\]

\begin{dinhly}[Trọng số công khai bảo toàn đảm bảo]
\label{thm:trem}
Với mọi hàm trọng số $w: V \to (0, 1]$ không phụ thuộc vị trí thật $x_t$ (chỉ
phụ thuộc $z_{<t}$, $\Delta t$ và các hằng số công khai), cơ chế T-REM thỏa
$\eps$-Geo-I tại mỗi bước: với mọi $x, x'$ và $v \in V$,
\[
\Pr[\mathrm{TREM}(x) = v \mid z_{t-1}] \le e^{\eps\,\dEu(x,x')}\,
\Pr[\mathrm{TREM}(x') = v \mid z_{t-1}].
\]
\end{dinhly}

\begin{proof}
Lặp lại chứng minh Định lý~\ref{thm:rem} với $Z_w(x) = \sum_u w(u)
e^{-\frac{\eps}{2}\dEu(x,u)}$. Thừa số $(A)$ không đổi vì $w(v)$ xuất hiện ở cả
tử và mẫu nên triệt tiêu:
$(A) = e^{\frac{\eps}{2}(\dEu(x',v)-\dEu(x,v))} \le e^{\frac{\eps}{2}\dEu(x,x')}$.
Với $(B)$: từng số hạng thỏa $w(u) e^{-\frac{\eps}{2}\dEu(x',u)} \le
e^{\frac{\eps}{2}\dEu(x,x')} w(u) e^{-\frac{\eps}{2}\dEu(x,u)}$ (trọng số $w(u)
> 0$ giữ nguyên hai vế), nên $Z_w(x') \le e^{\frac{\eps}{2}\dEu(x,x')} Z_w(x)$.
Kết luận như Định lý~\ref{thm:rem}. Điều kiện then chốt là $w$ đo được theo
thông tin công khai: nếu $w$ phụ thuộc $x_t$, thừa số $(A)$ không triệt tiêu và
lập luận sụp đổ --- đây chính là lỗi của bước làm mượt có kiểm tra QoS trong
baseline (Mệnh đề~\ref{prop:cap}).
\end{proof}

\begin{nhanxet}[Ngân sách toàn quỹ đạo]
Chuỗi $T$ lần công bố nhiễu độc lập tiêu tốn $\eps_{\mathrm{traj}} = T\eps$ theo
kết hợp tuần tự thích nghi (trọng số $w_t$ chọn theo đầu ra trước đó là hợp lệ
trong khuôn khổ adaptive composition). Đây là trần worst-case trung thực. SM-REM
(Mục~\ref{sec:smrem}) \emph{không} siết được trần này về mức toàn-quỹ-đạo: nó chỉ
đạt đảm bảo cho \emph{vị trí tĩnh lặp lại} (Định lý~\ref{thm:smrem}), còn pattern
revisit thì rò rỉ (Nhận xét~\ref{rmk:notrajectory}).
\end{nhanxet}

\section{SM-REM: Memoization chống tấn công averaging tại stay-point}
\label{sec:smrem}

Mọi cơ chế nhiễu độc lập theo thời gian --- kể cả REM/T-REM --- đều rò rỉ dưới
\emph{báo cáo lặp từ một vị trí tĩnh}: kẻ tấn công thấy $n$ release độc lập
$z_1,\dots,z_n$ của cùng một điểm thật $x$ (ví dụ ở nhà qua đêm) và lấy trung
bình; ước lượng hội tụ về $x$ với sai số $O(1/\sqrt{n})$.
Đây là phát biểu hình thức trong~\cite{chatzikokolakis2014predictive} (``khi
người dùng không di chuyển\ldots các vị trí báo cáo tập trung quanh vị trí thật,
lộ hoàn toàn nó khi số truy vấn tăng''), và là cơ chế đằng sau vụ Strava khôi
phục địa chỉ nhà (84--95\%, \cite{hassan2018run,dhondt2022run}) cùng home
fingerprinting của data broker. Đây chính là kịch bản S4 mà bản PTPPM lưới đường
mới nhất~\cite{min2025road} bỏ ngỏ: vì cơ chế của họ là \emph{stateless}
(Permute-and-Flip lấy mẫu độc lập mỗi timestamp), báo cáo lặp cùng một vị trí thật
vẫn bị averaging --- đây là suy luận của luận văn từ thiết kế của họ, không phải
một trích dẫn thừa nhận.

\paragraph{Cơ chế.} SM-REM phân hoạch không gian bằng một \emph{lưới công khai cố
định} kích thước $g$, với đại diện công khai $\mathrm{rep}(C)$ (tâm ô) mỗi ô. Lần
đầu điểm thật rơi vào ô $C$, một release $r_C$ được lấy mẫu bằng luật T-REM
\emph{tại $\mathrm{rep}(C)$} (không tại toạ độ thật) và lưu lại; mọi báo cáo sau có
điểm thật rơi vào $C$ đều trả về \emph{cùng} $r_C$. Sample tại $\mathrm{rep}(C)$
làm phân phối release của mọi điểm trong một ô giống hệt nhau --- điều kiện cho
phát biểu mức-ô. Quyết định memoize chỉ dựa trên lưới công khai (memoize mọi lần),
không có phép thử phụ thuộc dữ liệu.

\begin{dinhly}[Đảm bảo cho vị trí tĩnh lặp lại --- phạm vi hẹp có chủ ý]
\label{thm:smrem}
Cố định một vị trí thật $x$ thuộc ô $C$ và cho người dùng báo cáo $x$ đúng $T$
lần. Gọi $Z \sim \mathrm{REM}(\mathrm{rep}(C))$ là release đầu. Transcript
$(Z,\dots,Z)$ là hậu xử lý xác định của \emph{một} mẫu $Z$, nên: (i) thừa hưởng
đảm bảo $\eps$-Geo-I mức-ô của $Z$; và (ii) mang đúng lượng thông tin của một
release bất kể $T$ --- trung bình số học của $T$ bản sao giống hệt không giảm
phương sai, nên tấn công averaging (S4) không thu được gì.
\end{dinhly}

\begin{proof}
Sau khi $Z$ được lấy mẫu, mọi phát ra sau là tra bảng xác định trả về $Z$ --- một
hàm của riêng $Z$ --- nên bất biến hậu xử lý (\cite{takagi2019geo}, bổ đề hậu xử
lý) cho (i). Với (ii), phân phối của $(Z,\dots,Z)$ đồng nhất với phân phối của
$Z$; ước lượng trung bình $\frac1T\sum Z = Z$ có cùng phân bố với một mẫu đơn, độc
lập với $T$.
\end{proof}

\begin{nhanxet}[Điều KHÔNG được khẳng định --- và vì sao (V-003)]
\label{rmk:notrajectory}
Ta \emph{không} khẳng định một định lý toàn-quỹ-đạo kiểu
``$\eps\cdot(\text{số ô riêng biệt})$-Geo-I trên quỹ đạo tùy ý''. Memoization
chính xác làm pattern hit/miss của cache trở thành hàm xác định của \emph{pattern
revisit bí mật}, và pattern này rò rỉ: với hai quỹ đạo $X=(a,a)$ và $X'=(a,b)$
($b$ ở ô khác), biến cố $\{z_2 \ne z_1\}$ có xác suất $0$ dưới $X$ nhưng $>0$ dưới
$X'$ --- likelihood ratio vô hạn một chiều. Lưới công khai không làm một lần
revisit bí mật trở thành công khai. Một đảm bảo bảo vệ cả pattern revisit đòi hỏi
\emph{quyết định reuse tự nó phải riêng tư vi phân} (phép thử private kiểu
predictive mechanism~\cite{chatzikokolakis2014predictive}) --- được đặt làm hướng
phát triển (Chương~\ref{ch:tongket}), không phải phần đã chứng minh ở đây.
\end{nhanxet}

\begin{nhanxet}[Cũng không phải Geo-I Euclid --- gián đoạn ở biên]
\label{rmk:boundary}
Ngay ở mức một release, quantize $x \to \mathrm{rep}(C)$ là rời-rạc-hóa phía đầu
vào: hai điểm sát nhau hai bên biên ô nhận release độc lập, tỷ số tới
$\exp(\eps\cdot\text{cellwidth})$. SM-REM vì thế \emph{không} thỏa $\eps$-Geo-I
Euclid như REM/T-REM --- \textbf{không thể vừa có Geo-I Euclid sạch vừa có
memoization chống averaging từ cùng một tham số}; đây đúng là lý do Android
\texttt{LocationFudger} ghép snap-to-grid với offset ngẫu nhiên bền. Mặc định
$g=60$m; release snap về đỉnh đường dùng chung nên giá trị nhất quán không tự trở
thành fingerprint duy nhất.
\end{nhanxet}

\section{PR-SM-REM: reuse riêng tư và đảm bảo w-event}
\label{sec:prsmrem}

Nhận xét~\ref{rmk:notrajectory} chỉ ra vì sao SM-REM \emph{không} có đảm bảo toàn
quỹ đạo: quyết định reuse (hit/miss cache) là hàm xác định của pattern revisit bí
mật. Cách vá nguyên lý (predictive mechanism~\cite{chatzikokolakis2014predictive};
kế toán w-event~\cite{kellaris2014wevent}) là làm chính \emph{quyết định reuse}
riêng tư vi phân.

\paragraph{Cơ chế.} Ở mỗi bước $t$, dự đoán $\tilde z_t$ = \emph{điểm đã công bố
trước đó} $z_{t-1}$ (parrot predictor --- công khai, chỉ là hàm của transcript).
Chạy \emph{phép thử ngưỡng có nhiễu}:
\[
d(x_t, \tilde z_t) + \mathrm{Lap}(1/\eps_{\text{test}}) \le \theta
\quad\Rightarrow\quad \text{REUSE } \tilde z_t;\qquad
\text{ngược lại} \Rightarrow \text{RESAMPLE bằng REM}.
\]
Reuse tốn đúng $\eps_{\text{test}}$; resample tốn $\eps_{\text{test}}+\eps$.
$\theta$ là bán kính reuse công khai; trạng thái chỉ phụ thuộc thông tin công
khai và các release quá khứ, không bao giờ phụ thuộc $x_{\le t}$ thô.

\begin{dinhly}[Chặn per-step của PR-SM-REM]
\label{thm:prsmrem}
Kernel một bước là hỗn hợp hai nhánh: tại $z=h$ có khối lượng
$q_x(h)+(1-q_x(h))R_x(h)$; tại $z\ne h$ có $(1-q_x(h))R_x(z)$, với
$q_x(h)=\Pr[d(x,h)+\mathrm{Lap}(1/\eps_{\text{test}})\le\theta]$ và $R$ là kernel
REM với ngân sách $\eps_{\text{release}}$. Vì thống kê $d(x,h)$ có độ nhạy metric
bằng 1 theo $x$ và $h$ công khai, nhánh RESAMPLE mang \emph{cả} thừa số phép thử
lẫn release REM, nên chặn per-step worst-case là
\[
\frac{K_x(z\mid h)}{K_{x'}(z\mid h)} \le \exp\!\big((\eps_{\text{test}}+\eps_{\text{release}})\,d(x,x')\big),
\]
tức \textbf{$(\eps_{\text{test}}+\eps_{\text{release}})$-Geo-I} (bước đầu tiên,
chưa có dự đoán, chỉ tốn $\eps_{\text{release}}$). Do đó biến cố $\{z\ne h\}$ có
tỷ số \emph{hữu hạn} $\le\exp((\eps_{\text{test}}+\eps_{\text{release}})d(x,x'))$
--- đây mới là điều đóng lỗ hổng tỷ-số-vô-hạn của memo chính xác
(Nhận xét~\ref{rmk:notrajectory}), \emph{không} phải chặn chỉ-$\eps_{\text{test}}$.
\end{dinhly}

\begin{proof}
Điều kiện trên \emph{cùng một public history} (nên $h$, $q$'s form giống nhau ở
hai input). Đặt $d=d(x,x')$ và $c=e^{(\eps_{\text{test}}+\eps_{\text{release}})d}$.
Ba chặn thành phần: $q_x\le e^{\eps_{\text{test}}d}q_{x'}$ và
$1-q_x\le e^{\eps_{\text{test}}d}(1-q_{x'})$ (query $d(\cdot,h)$ độ nhạy 1 + nhiễu
Laplace), và $R_x(z)\le e^{\eps_{\text{release}}d}R_{x'}(z)$ (Định lý~\ref{thm:rem}).
\emph{Nhánh $z\ne h$:} $K_x(z\mid h)=(1-q_x)R_x(z)\le e^{\eps_{\text{test}}d}
e^{\eps_{\text{release}}d}(1-q_{x'})R_{x'}(z)=c\,K_{x'}(z\mid h)$.
\emph{Nhánh $z=h$:} \emph{cả hai} số hạng chịu cùng hệ số $c$ (số hạng $q_x$ chịu
$e^{\eps_{\text{test}}d}\le c$; số hạng $(1-q_x)R_x(h)$ chịu $c$), nên tổng
$q_x+(1-q_x)R_x(h)\le c\,[q_{x'}+(1-q_{x'})R_{x'}(h)]=c\,K_{x'}(h\mid h)$ ---
không dùng ``tỷ số của tổng'' một cách không giải thích. Liệt kê hai-vertex
($a{=}0$m, $b{=}100$m, $h{=}a$, $\theta{=}200$m,
$\eps_{\text{test}}{=}\eps_{\text{release}}{=}0{,}01$; $z{=}b$): tỷ số $=e^{1.5}$,
vượt chặn cũ $e^{1}$ nhưng thỏa $e^{2}$ --- được test hồi quy hóa ở
\texttt{tests/test\_pr\_privacy\_bound.py}.
\end{proof}

\begin{nhanxet}[Ideal kernel vs hiện thực dấu-phẩy-động]
\label{rmk:finiteprec}
Mọi chặn ở trên là cho \emph{ideal real-arithmetic kernel}. Sampler Gumbel-max và
phép thử Laplace trong hiện thực dùng RNG 53-bit hữu hạn nên có span bị chặn: một
đỉnh có logit thấp hơn leader quá $\approx$40,34 \emph{không bao giờ} được chọn,
tạo zero-support phụ thuộc input (tỷ số vô hạn ở mức executable). Vì vậy luận văn
\emph{không} tuyên bố hiện thực float thỏa pure $\eps$-Geo-I; nó là một xấp xỉ số
học của ideal kernel. Chứng minh pure-DP hữu-hạn-độ-chính-xác (exact discrete
sampler với tinh chỉnh bit không giới hạn) là hướng phát triển. Xem
\texttt{tests/test\_sampler\_support.py} (verifier R2-006).
\end{nhanxet}

\begin{nhanxet}[Từ per-step tới quỹ đạo --- và điều chưa cài (R2-002)]
\label{rmk:wevent}
Kết hợp tuần tự trên cửa sổ $w$ (metric $D_\infty(X,X')=\max_t d(x_t,x'_t)$) cho
worst-case $w(\eps_{\text{test}}+\eps_{\text{release}})$ mỗi cửa sổ (cửa sổ đầu:
$\eps_{\text{release}}+(w-1)(\eps_{\text{test}}+\eps_{\text{release}})$). Hiện
thực \emph{chưa} có budget manager/privacy-filter theo cửa sổ; \texttt{n\_resample}
chỉ là chẩn đoán hậu nghiệm, \emph{không} phải một guarantee ex-ante. Muốn tuyên
bố một scalar $\eps_w$ nhỏ hơn worst-case cần thêm một allocator theo cửa sổ và
lấy supremum trên mọi public transcript --- đặt làm hướng phát triển
(Chương~\ref{ch:tongket}). Để so sánh công bằng với REM $\eps$-Geo-I, đặt
$\eps_{\text{test}}+\eps_{\text{release}}=\eps$ (thực nghiệm dùng
$\eps_{\text{test}}=\eps_{\text{release}}=\eps/2$).
\end{nhanxet}

\begin{nhanxet}[Đánh đổi thực nghiệm và phạm vi]
So với SM-REM (memo chính xác), PR-SM-REM \emph{yếu hơn} về chống-averaging thực
nghiệm nhưng có chặn per-step hữu hạn hợp lệ. Trạng thái chỉ là \emph{release
trước đó} nên nó là \emph{predictive reuse} của điểm liền trước, \emph{không phải}
memoization bền vững: mẫu rời-rồi-quay-lại (nhà$\to$chỗ làm$\to$nhà) không được
nhớ (verifier R2-007). $\theta$ phải công khai.
\end{nhanxet}

\section{Thuật toán và độ phức tạp}

\begin{enumerate}
  \item \textbf{Tiền xử lý} (một lần cho vùng quan tâm): dựng đồ thị đường OSM,
  chiếu các đỉnh về hệ tọa độ phẳng cục bộ, xây KD-tree trên $V$.
  \item \textbf{Mỗi điểm thật $x_t$}: tính logits
  $-\tfrac{\eps}{2}\dEu(x_t, u) + \log w_t(u)$ trên \emph{toàn bộ} tập đỉnh cố
  định $u \in V$ (REM: $w_t \equiv 1$) --- không cắt ngưỡng quanh $x_t$
  (Nhận xét~\ref{rmk:cutoff}); lấy mẫu bằng Gumbel-max trong miền log; công bố
  đỉnh được chọn; cập nhật $z_t$.
\end{enumerate}

Chi phí mỗi điểm là $O(|V|)$ phép tính vector hóa (một 2-norm trên $|V|$ đỉnh +
một argmax); với đồ thị Bắc Kinh $|V|=13\,813$ đỉnh dùng trong thực nghiệm, chi
phí này dưới mili-giây trong reference implementation Python. (Đây là chi phí của
riêng bước lấy mẫu; con số benchmark tổng gồm cả metrics và attackers nên không
dùng làm latency claim on-device --- cần đo cách ly trên phần cứng mobile, xem
Chương~\ref{ch:tongket}.)
